\scanpage{50}
\diagram{\begin{tikzcd}[ampersand replacement=\&,row sep=small] (\deleted{J}B,\sigma)\arrow[d]\&[-9pt] \deleted{J}B\ \text{ens. fini, }\sigma\ \text{involution dans }\deleted{J}B\ \text{sans pt fixe}\\ E=E(\deleted{J}B,\sigma)=\operatorname{Hom}_{\{\pm1\}}(\deleted{J}B,\mathbb R)\& \end{tikzcd}}
\[C=C(\deleted{J}B,\sigma)=\{\varphi\in E\mid\lvert\varphi(\deleted{j}b)\rvert\leq1\quad\forall \deleted{j}b\in\deleted{J}B\}\qquad\deleted{J}B\hookrightarrow C\hookrightarrow E\]

\uline{Proposition} Pour toute section partielle \(I'\subset \deleted{J}B\) de \(\deleted{J}B\) sur \(I=\deleted{J}B/\sigma\), soit
\[F(I')=\left\{\begin{array}{l|l}\varphi\in C&\varphi(j)=1\quad\forall j\in I'\\[-2pt]\deleted{\unsure{\varphi(j)}}&\inserted{\varphi(j)\neq\pm1\quad\forall j\in \deleted{J}B-I'-\sigma I'}\end{array}\right\}\deleted{\unsure{ou}}.\]
Alors l'application \(I'\longmapsto F(I')\) est une application bijective, renversant l'ordre
\[\operatorname{Spd}(J/I)\xrightarrow{\sim}(\operatorname{Fac}\deleted{\unsure{f}}\,(C)-\{0\})\]
\deleted{\unsure{faces|fermées}}

Ceci précise de façon simple le \deleted{\unsure{le}} polyèdre combinatoire associé au polyèdre \(C=C(J,\sigma)\) en termes de \((J,\sigma)\), comme étant (modulo l'adjonction d'un plus petit élément) l'ens. ordonné opposé de l'ens des parties \(I'\) de \(J\) telles que \(I'\cap\sigma I'=\varnothing\).

% margin note written sideways
\margin{\uline{NB} Dans le cas du cube (\(p\) de degré \(2\)) les él. de \(B\) correspondent aux facettes de \unsure{codim un|codim \(1\)}, et pour une telle \(F\), \(\sigma F\) est l'unique facette \inserted{\unsure{de dim \(n-1\)|de codim \(1\)}} \deleted{\unsure{qui}} qui ne rencontre pas \(F\) (en fait, les \uline{hyperpl. d'appui correspondants sont parallèles})}

\uline{Proposition} Soit \(\deleted{J}B\xrightarrow{p} I\) une application surjective, et soit \(P(p)\) l'ens ordonné par inclusion des sections partiellement définies de \(\deleted{J}B\) sur \(I\). Alors \(p\longmapsto P(p)\) est un foncteur pl. fidèle du groupoïde des surjections (avec comme \deleted{\unsure{ses}} morphismes les isom.
\diagram{\begin{tikzcd}[ampersand replacement=\&,row sep=small,column sep=small]\deleted{J}B\arrow[r,"\sim"]\arrow[d,"p"']\&\deleted{J}B'\arrow[d,"p'"]\\I\arrow[r,"\sim"]\&I'\end{tikzcd}}
) dans la catégorie des ens. ordonnés

Car on reconstitue \(B\hookrightarrow P\), ens des \unsure{élts|épts|pts} \deleted{\unsure{maxi}}minimaux de \(P-\unsure{0|\varnothing}\)\deleted{\unsure{\(\}\)}} \inserted{\(\hookrightarrow\) plus petit élt de \(P\)} et la relation d'équivalence définie par \(p\) par
\[p(b)=p(b')\Longleftrightarrow(b=b'\ \text{ou}\ C_b\cap C_{b'}=0),\quad\text{où }\forall b\in B,\]
\(C_b\) est \deleted{l'ens.} \inserted{sous-ens. de l'ens \(S\)} (\deleted{des éléments} maxima \(c\) de \(P\) (\(S\sim\prod_{i\in I} B_i\)) \deleted{tels} qui majorent \(b\).
% archive sheet number: 49
